\documentclass{article}
\usepackage{graphicx} % Required for inserting images

\begin{document}
\begin{titlepage}
\centering

 \includegraphics[width=0.7\textwidth]{thapar_logo.jpg}
\rule{0pt}{3cm}

\vspace{1cm}

{\Large \textbf{THAPAR INSTITUTE OF ENGINEERING AND TECHNOLOGY}}\\
{\large Department of Computer Science and Engineering}

\vspace{1.5cm}

{\Huge \textbf{Project Proposal}}\\[0.5cm]
{\LARGE \textbf{Offline Hand Sign to Text Recognition System for American Sign Language}}

\vspace{1.5cm}

\textbf{Course Name \& Code:} UCS503 -- Software Engineering\\
\textbf{Degree \& Year:} B.Tech (3rd Year)\\
\textbf{Institute:} Thapar Institute of Engineering and Technology\\[0.5cm]

\textbf{Submitted By:}\\
Abhishek Srivastava \hfill - 1024030236\\
Keshav Aggarwal \hfill - 1024030234\\
Harshine Singh \hfill - 1024030232\\
Bhuvaneesh Singh Grewal \hfill - 1024030235\\[1cm]


\textbf{Academic Year:} 2026--27 ODD SEMESTER

\vfill

\end{titlepage}

\maketitle

\section{Introduction and Overview}

\subsection{Purpose}
The purpose of \textit{Coffee With Bro} is to develop an offline American Sign Language (ASL) hand-sign-to-text recognition system. The system is intended to recognize static ASL alphabet letters and numerical digits from live webcam input and display the corresponding characters as text.

\subsection{Scope}
The system will recognize the 26 ASL alphabet letters and 10 numerical digits, giving a total of 36 classes. It will capture hand-sign images through a webcam, process the captured frames, classify them using a CNN trained from scratch, and display the predicted character. The system will operate without cloud-based inference services, third-party APIs, pretrained models, or internet connectivity during inference.

\subsection{Definitions and Acronyms}
\begin{itemize}
    \item \textbf{ASL}: American Sign Language.
    \item \textbf{CNN}: Convolutional Neural Network.
    \item \textbf{DL4J}: DeepLearning4j, a Java-based deep learning framework.
    \item \textbf{API}: Application Programming Interface.
    \item \textbf{NFR}: Non-Functional Requirement.
    \item \textbf{UI}: User Interface.
\end{itemize}


\section{Overall Description}

\subsection{Product Perspective}
\textit{Coffee With Bro} is a browser-based sign recognition system consisting of three primary components: a frontend for webcam access and result display, a Spring Boot backend for serving predictions, and a CNN trained using DeepLearning4j for sign classification. The frontend captures webcam frames and communicates with the backend, which processes the frames using the trained model and returns the predicted character.

\subsection{User Characteristics}
The system is designed for users who need or wish to convert static ASL hand signs into readable text. Users are expected to have basic familiarity with using a computer, a web browser, and a webcam. No specialized technical knowledge is required to operate the system.

\subsection{Constraints and Assumptions}
\begin{itemize}
    \item The system must operate locally without internet connectivity during inference.
    \item No pretrained models or transfer learning will be used.
    \item The CNN will be trained from scratch using the project's training pipeline.
    \item The system requires a compatible webcam and local computing hardware.
    \item The current system is limited to static ASL alphabet letters and numerical digits.
    \item The dataset is maintained separately from the source-code repository.
\end{itemize}


\section{Functional Requirements}

\subsection{Features and Actions}
The system shall:
\begin{enumerate}
    \item Access the user's webcam through the browser.
    \item Capture image frames containing hand signs.
    \item Send captured frames to the backend for classification.
    \item Process the input using the trained CNN model.
    \item Classify the input into one of 36 ASL classes.
    \item Return the predicted letter or digit to the frontend.
    \item Display the predicted character to the user.
    \item Accumulate recognized characters into a running text buffer in subsequent development.
\end{enumerate}

\subsection{Use Cases}
The primary use case is ASL sign recognition:
\begin{enumerate}
    \item The user opens the web application.
    \item The user grants access to the webcam.
    \item The user presents an ASL hand sign to the camera.
    \item The frontend captures a frame.
    \item The frame is sent to the Spring Boot backend.
    \item The trained CNN processes the frame.
    \item The backend returns the predicted ASL character.
    \item The frontend displays the prediction.
\end{enumerate}

\subsection{Data Requirements}
The system uses the ASL-HG dataset containing 36,000 images across 36 classes: A--Z and 0--9. The raw images are resized to $128 \times 128$ pixels, shuffled, and divided into an 80/20 training and testing split using a fixed random seed. The trained CNN model is then used to classify live webcam frames.


\section{Non-Functional Requirements (NFRs)}

\subsection{Performance}
\begin{itemize}
    \item The CNN shall target at least 90\% top-1 classification accuracy on the held-out test set.
    \item The system shall provide sufficiently low inference latency for interactive webcam-based use.
    \item Model performance shall be evaluated using held-out test data after training.
\end{itemize}

\subsection{Security}
\begin{itemize}
    \item The system shall not depend on third-party cloud inference services.
    \item No API keys shall be required for model inference.
    \item Inference shall be performed locally on the user's hardware.
\end{itemize}

\subsection{Reliability and Usability}
\begin{itemize}
    \item The preprocessing pipeline shall produce a reproducible train/test split using a fixed random seed.
    \item The frontend, backend, and preprocessing components shall be independently testable.
    \item The user interface shall provide clear visual feedback for recognized characters.
    \item The system shall support live webcam-based interaction at an acceptable interactive pace.
\end{itemize}


\section{External Interfaces}

\subsection{User Interfaces}
The system shall provide a browser-based interface that:
\begin{itemize}
    \item Provides access to the user's webcam.
    \item Displays the live camera input.
    \item Displays the predicted ASL letter or digit.
    \item Provides a running text output as characters are recognized.
\end{itemize}

\subsection{Software and Hardware Interfaces}
The system interfaces with:
\begin{itemize}
    \item A webcam through the browser's MediaDevices/getUserMedia API.
    \item A Spring Boot backend for frame classification.
    \item A DeepLearning4j CNN for machine-learning inference.
    \item Local computer hardware for model training and inference.
    \item The browser environment for frontend execution.
\end{itemize}
\end{document}